% HISTORICAL (resolved 2026-07-10): this counterexample concerns the June 2026
% updated_cliff.tex gates. The July 1, 2026 universal_new_gates.tex replaces the
% gates with an orthogonal-axes pair, adopting the second repair suggested below.
\documentclass[11pt]{article}

\usepackage[margin=1in]{geometry}
\usepackage{amsmath}
\usepackage{amssymb}

\newcommand{\Tr}{\operatorname{Tr}}
\newcommand{\ii}{i}
\newcommand{\Rz}{R_z}
\newcommand{\Gone}{G_1}
\newcommand{\Gtwo}{G_2}
\newcommand{\pauli}{\sigma}

\begin{document}

\section*{A counterexample to the three-factor $G_1,G_2$ Euler step}

Let
\[
  \Gone = e^{-\ii\pi/4}THTH,
  \qquad
  \Gtwo = e^{-\ii\pi/4}HTHT.
\]
The claim under consideration is that every determinant-one one-qubit unitary
is, up to global phase, of the form
\[
  \Gone^a\Gtwo^b\Gone^c
\]
for some real numbers $a,b,c$.  The following exact calculation shows that this
claim is false for these particular gates.

Write a determinant-one one-qubit unitary in the form
\[
  X=tI+\ii\,w\cdot\pauli,
  \qquad
  t\in\mathbb{R},\quad w\in\mathbb{R}^3,\quad t^2+\|w\|^2=1.
\]
Also write
\[
  \Gone=\cos(\alpha)I+\ii\sin(\alpha)n_1\cdot\pauli,
  \qquad
  \Gtwo=\cos(\alpha)I+\ii\sin(\alpha)n_2\cdot\pauli,
\]
where $n_1,n_2$ are unit vectors.  The two gates have the same trace, so the
same $\alpha$ can be used, up to sign choices for the axes.  These sign choices
do not affect the argument below, since only squares occur.

Direct matrix calculation gives
\[
  \Tr(\Gone)=\Tr(\Gtwo)=1+\frac{1}{\sqrt2}.
\]
Thus
\[
  \cos(\alpha)
  =
  \frac12+\frac{1}{2\sqrt2}
  =:c,
  \qquad
  \sin^2(\alpha)=1-c^2=\frac{5-2\sqrt2}{8}.
\]
From the exact matrix for $\Gone$,
\[
  (\Gone)_{00}
  =
  \frac{(1-\ii)(1+\sqrt2+\ii)}{4}
  =
  c-\frac{\ii}{2\sqrt2}.
\]
Since the $(0,0)$ entry of
$\cos(\alpha)I+\ii\sin(\alpha)n_1\cdot\pauli$ is
$\cos(\alpha)+\ii\sin(\alpha)(n_1)_z$, we get
\[
  \sin^2(\alpha)(n_1)_z^2=\frac18,
  \qquad
  (n_1)_z^2=\frac{1}{5-2\sqrt2}.
\]

Let
\[
  k=\langle n_1,n_2\rangle .
\]
Another direct matrix calculation gives
\[
  \Gone\Gtwo
  =
  \frac12
  \begin{pmatrix}
    1-\ii & -\ii\sqrt2 \\
    -\ii\sqrt2 & 1+\ii
  \end{pmatrix},
\]
so $\Tr(\Gone\Gtwo)=1$.  On the other hand,
\[
  \frac12\Tr(\Gone\Gtwo)
  =
  \cos^2(\alpha)-\sin^2(\alpha)k.
\]
Therefore
\[
  \frac12=c^2-\sin^2(\alpha)k,
  \qquad
  k=\frac{2\sqrt2-1}{5-2\sqrt2}.
\]
Now
\[
\begin{aligned}
  k^2-(n_1)_z^2
  &=
  \frac{(2\sqrt2-1)^2}{(5-2\sqrt2)^2}
  -
  \frac{1}{5-2\sqrt2}                                      \\
  &=
  \frac{(2\sqrt2-1)^2-(5-2\sqrt2)}{(5-2\sqrt2)^2}             \\
  &=
  \frac{4-2\sqrt2}{(5-2\sqrt2)^2}
  >0.
\end{aligned}
\]
Hence
\[
  (n_1)_z^2<k^2.
\]

Next, every three-factor product about these two axes satisfies a necessary
constraint.  Let
\[
  R_n(x)=\cos(x)I+\ii\sin(x)n\cdot\pauli
\]
and suppose
\[
  R_{n_1}(a)R_{n_2}(b)R_{n_1}(c)=tI+\ii w\cdot\pauli.
\]
Diagonalize $n_1\cdot\pauli$.  In that basis, the $(+1,+1)$ entry is
\[
  t+\ii\langle w,n_1\rangle
  =
  e^{\ii(a+c)}\bigl(\cos(b)+\ii k\sin(b)\bigr).
\]
Taking squared absolute values gives
\[
  t^2+\langle w,n_1\rangle^2
  =
  \cos^2(b)+k^2\sin^2(b)
  \geq k^2.
\]

Now take the exact target
\[
  \Rz(\pi)
  =
  \begin{pmatrix}
    -\ii & 0\\
    0 & \ii
  \end{pmatrix}
  =
  -\ii Z.
\]
In the notation above, this has $t=0$ and $w=-e_z$.  Therefore
\[
  t^2+\langle w,n_1\rangle^2=(n_1)_z^2<k^2,
\]
which violates the necessary constraint.  So there are no real numbers
$a,b,c$ such that
\[
  \Rz(\pi)=\Gone^a\Gtwo^b\Gone^c.
\]
The same conclusion holds up to global phase.  Both sides have determinant
one, so a possible global phase must square to one and is therefore $\pm1$.
Multiplying by $-1$ changes $(t,w)$ to $(-t,-w)$, and the quantity
$t^2+\langle w,n_1\rangle^2$ is unchanged.

Thus the exact gate $\Rz(\pi)$ is a counterexample to the claimed global
three-factor decomposition for these $G_1,G_2$ axes.

\section*{A corrected half-angle repair}

The global three-factor statement is false, but the proof can be repaired for
the gates that are actually needed, namely the $z$-rotations.  The replacement
statement is local: sufficiently small $z$-rotations have the desired
three-factor form.

The criterion above says that a target
\[
  X=tI+\ii w\cdot\pauli
\]
can have the form
\[
  R_{n_1}(a)R_{n_2}(b)R_{n_1}(c)
\]
only if
\[
  t^2+\langle w,n_1\rangle^2\geq k^2.
\]
Conversely, this same inequality is sufficient.  To see this, diagonalize
$n_1\cdot\pauli$.  In that basis, the top-left entry of
$R_{n_1}(a)R_{n_2}(b)R_{n_1}(c)$ is
\[
  e^{\ii(a+c)}(\cos(b)+\ii k\sin(b)).
\]
Thus the squared norm of the top-left entry ranges from $k^2$ to $1$ as $b$
varies.  Once $b$ is chosen to match this norm, the two remaining phases are
matched by choosing $a+c$ and $a-c$.  This gives the needed $a$ and $c$.

Now consider a $z$-rotation
\[
  \Rz(\varphi)=\cos(\varphi/2)I-\ii\sin(\varphi/2)Z.
\]
For this gate,
\[
  t^2+\langle w,n_1\rangle^2
  =
  \cos^2(\varphi/2)+(n_1)_z^2\sin^2(\varphi/2).
\]
If $|\varphi|\leq \pi/2$, then
\[
  \sin^2(\varphi/2)\leq \frac12,
\]
and therefore
\[
  \cos^2(\varphi/2)+(n_1)_z^2\sin^2(\varphi/2)
  \geq
  \frac{1+(n_1)_z^2}{2}.
\]
For the specific $G_1,G_2$ axes computed above,
\[
  \frac{1+(n_1)_z^2}{2}-k^2
  =
  \frac{50-31\sqrt2}{289}
  >0.
\]
Hence every $z$-rotation with $|\varphi|\leq\pi/2$ has a three-factor form
\[
  \Rz(\varphi)\sim \Gone^a\Gtwo^b\Gone^c
\]
for suitable real numbers $a,b,c$.

Finally, let $\theta$ be arbitrary.  After reducing $\theta$ modulo $2\pi$ up
to global phase, assume $|\theta|\leq\pi$.  Put $\varphi=\theta/2$.  Then
$|\varphi|\leq\pi/2$, so the preceding paragraph gives
\[
  \Rz(\varphi)\sim \Gone^a\Gtwo^b\Gone^c.
\]
Since
\[
  \Rz(\theta)=\Rz(\varphi)^2,
\]
we obtain the corrected real-power decomposition
\[
  \Rz(\theta)
  \sim
  (\Gone^a\Gtwo^b\Gone^c)(\Gone^a\Gtwo^b\Gone^c).
\]
Thus the false three-factor statement can be replaced, for the $R_z$ family,
by a six-factor statement obtained from two half-rotations.  The rest of the
approximation argument can then proceed as before: approximate each real power
of $G_1$ and $G_2$ by an integer power, and use the product-distance lemma to
add the errors.

An alternative repair is to use an orthogonal-axis version of the
Boykin/Nielsen--Chuang argument.  With orthogonal axes, the three-factor Euler
decomposition is globally valid, so the counterexample above does not apply.

\end{document}
